%% -*- coding: utf-8 -*-
\documentclass[12pt,pagesize,paper=192mm:108mm,landscape]{scrbook} 
%1920x1080 1280x720
\areaset[current]{192mm}{108mm}
\usepackage{calc}
\usepackage[T2A]{fontenc}
\usepackage[utf8]{inputenc}
\usepackage[english,russian]{babel}
\usepackage{microtype}
\usepackage{misccorr}
\usepackage{cmap}
%\usepackage[unicode=true]{hyperref}
\usepackage{graphicx}
\usepackage{amssymb}
\usepackage{amsmath}
%\usepackage{srcltx}
\usepackage{textcomp}
\usepackage{xspace}
%научные символы и смайлики \smiley \frownie
\usepackage{wasysym}
\usepackage{ccicons}
\begin{document}
\begin{titlepage}
  \vspace*{-0.5em}
  \begin{center}    
    % \hspace*{3em}
    % \begin{minipage}[t]{3em}
    %   \includegraphics[width=\textwidth]{../BINP-logo}
    % \end{minipage}\hfill
    % \begin{minipage}{0.23\linewidth}
    % \includegraphics[width=\textwidth]{../NSU-logo}
    % \end{minipage}
    % \hfill
    % \hspace*{6em}

    
    % Кафедра теоретической физики физического факультета НГУ
    % \medskip

    % \Large
    % Профессор Грабовский А.\,В.
    % \smallskip

    % \huge
    % \textbf{Общая теория относительности}
    % \smallskip

    % \Large
    % Лекция № 10
     \vfill

     \normalsize
     \begin{minipage}{0.8\linewidth}
       Получение уравнений Эйнштейна в лоренцевской
       калибровке. Решение уравнений Эйнштейна в лоренцевской
       калибровке с тензором энергии импульса и в
       пустоте. Демонстрация того, что скорость гравитационной волны
       равна скорости света. Получение закона дисперсии гравитационной
       волны. Нахождение ``+'' и ``х'' поляризаций гравитационной волны в
       поперечно бесследовой калибровке. Демонстрация невозможности
       обнаружить гравитационную волну в одной точке. Изменение
       расстояния (пространственного интервала) между двумя близкими
       точками при прохождении гравитационной волны с различными
       поляризациями. Демонстрация поперечности гравитационной
       волны. Демонстрации возможности перехода в лоренцевскую
       калибровку. Демонстрация возможности перехода в поперечно
       бесследовую калибровку. Доказательство того, что в эргосфере
       частица с отрицательной энергией имеет отрицательный момент
       импульса, то есть при падении на~черную дыру она уменьшает ее
       массу и момент. Нахождение пространственной метрики на
       вращающемся диске.
     \end{minipage}
    \vfill

    \normalsize \ccbysa\hspace{0.5em}  Новосибирск 2022
  \end{center}
\end{titlepage}
\end{document}
